\documentclass[aspectratio=169,10pt,english]{beamer}

\input{../../../_preambulo_beamer.tex}
\input{../../../_comandos_beamer.tex}

\selectlanguage{english}

\title{MA1001B - Statistical Modeling for Decision Making}
\subtitle{Lesson 37: Two Independent Means (Two-Sample t)}
\author{}
\institute{Tecnol\'ogico de Monterrey}
\date{}

\begin{document}

\begin{frame}[plain]
  \vspace{1.2cm}
  {\Large\bfseries MA1001B - Statistical Modeling for Decision Making\par}
  \vspace{0.7cm}
  {\large Lesson 37: Two Independent Means (Two-Sample t)\par}
  \vspace{0.7cm}
  {\color{TecRojo}\rule{\textwidth}{1.2pt}}
  \vspace{0.8cm}

  MA1001B Faculty\\[0.3cm]
  Term\\[0.3cm]
  Tecnol\'ogico de Monterrey
\end{frame}

\begin{frame}{The Two-Shift Mean Question}
  \begin{alertblock}{Guiding question}
    How do we compare mean pick times from two independent shifts when each
    group has its own $s$?
  \end{alertblock}

  \vspace{0.6em}
  By the end of this lesson, you can:
  \begin{itemize}
    \item keep two independent samples unpooled;
    \item compute a Welch $t$ statistic and Satterthwaite $df$;
    \item decide $H_0:\mu_A=\mu_B$ from a two-sided p-value;
    \item refuse a causal reading of an observational grouping.
  \end{itemize}

  \vfill
  \scriptsize All pick times and experience values used today are synthetic. No real company data are used.
\end{frame}

\begin{frame}{Two Independent Samples, Two Standard Deviations}
  \small
  A synthetic warehouse samples $n_A=30$ picks from Shift A and $n_B=32$
  picks from Shift B. The associates are not paired, and the sample
  standard deviations need not be equal.

  \[
    t=\frac{\bar{x}_A-\bar{x}_B}{\sqrt{s_A^2/n_A+s_B^2/n_B}}.
  \]

  \begin{exampleblock}{Stated procedure: Welch}
    $H_0:\mu_A=\mu_B$ versus $H_1:\mu_A\neq\mu_B$, $\alpha=0.05$. Degrees of
    freedom are Satterthwaite's, not $n_A+n_B-2$.
  \end{exampleblock}
\end{frame}

\begin{frame}[fragile]{Step 1: Group Summaries}
  \begin{columns}[T]
    \begin{column}{0.42\textwidth}
      \small
      \begin{lstlisting}
mean_a = times_a.mean()
s_a = times_a.std(ddof=1)
diff = mean_a - mean_b
      \end{lstlisting}

      \begin{block}{Verified output}
        $\bar{x}_A=48.100885$, $s_A=4.660494$\\
        $\bar{x}_B=52.826941$, $s_B=6.325509$\\
        $\bar{x}_A-\bar{x}_B=-4.726057$
      \end{block}
    \end{column}
    \begin{column}{0.58\textwidth}
      \centering
      \includegraphics[width=\textwidth]{../figures/two_sample_t_01_group_summaries.png}
      \vspace{0.2em}
      \scriptsize Independent-shift boxplots from Step 1
    \end{column}
  \end{columns}
\end{frame}

\begin{frame}{Welch $t$ Does Not Pool the Variances}
  \small
  \[
    SE=\sqrt{\frac{4.660494^2}{30}+\frac{6.325509^2}{32}}=1.405128,
  \]
  \[
    t=\frac{-4.726057}{1.405128}=-3.363436,\qquad df=56.900423.
  \]
  The two-sided p-value is $0.001382$. Because $|t|>t^*=2.002541$ and
  $p<\alpha$, the test rejects $H_0:\mu_A=\mu_B$.
\end{frame}

\begin{frame}[fragile]{Step 2: Welch Two-Sample Decision}
  \begin{columns}[T]
    \begin{column}{0.40\textwidth}
      \small
      \begin{lstlisting}
se = sqrt(s_a**2/n_a + s_b**2/n_b)
t = (mean_a - mean_b) / se
p = 2 * t_dist.sf(abs(t), df)
      \end{lstlisting}

      \begin{block}{Verified output}
        $t=-3.363436$, $df=56.900423$\\
        $p=0.001382$\\
        Decision: reject $H_0$
      \end{block}
    \end{column}
    \begin{column}{0.60\textwidth}
      \centering
      \includegraphics[width=\textwidth]{../figures/two_sample_t_02_welch_test.png}
      \vspace{0.2em}
      \scriptsize Welch $t$ tails from Step 2
    \end{column}
  \end{columns}
\end{frame}

\begin{frame}{Step 3: Observational Groups Are Not an Experiment}
  \begin{columns}[T]
    \begin{column}{0.42\textwidth}
      \small
      Shift A mean experience: $4.484559$ years.

      \vspace{0.3em}
      Shift B mean experience: $1.801265$ years.

      \vspace{0.4em}
      Associates were not randomly assigned to shifts.

      \vspace{0.5em}
      \textbf{Limit:} a significant Welch $p=0.001382$ compares means. It
      does not identify the shift as a cause. Lesson 38 studies matched
      before-after pairs instead.
    \end{column}
    \begin{column}{0.58\textwidth}
      \centering
      \includegraphics[width=\textwidth]{../figures/two_sample_t_03_observational_limit.png}
      \vspace{0.2em}
      \scriptsize Experience confounder from Step 3
    \end{column}
  \end{columns}
\end{frame}

\begin{frame}{Concept Check}
  \begin{enumerate}
    \item Why are these two shifts independent samples rather than pairs?
    \item Write the Welch standard error using $s_A$, $s_B$, $n_A$, and $n_B$.
    \item If $p=0.001382$ and $\alpha=0.05$, what is the decision for
      $H_0:\mu_A=\mu_B$?
    \item Why does a difference in mean experience block a causal claim?
  \end{enumerate}

  \vspace{0.6em}
  \begin{block}{Connection to Lesson 38}
    The session can continue directly into a paired $t$ test, where the same
    units are measured before and after a change.
  \end{block}

  \vfill
  \scriptsize Source: OpenStax, \textit{Introductory Business Statistics 2e},
  Chapter 10, Section 10.1. CC BY-NC-SA.
\end{frame}

\appendix



\end{document}
